\section{Overlap graph, sharp threshold, and branch locus}\label{sec:stab:conjugacy}

The local residue fibers from Section~\ref{sec:stabilization} become a dynamical question only after one imposes overlap. The theorem chain in this section is the structural core of the paper: Theorem~\ref{thm:finite-state-lift} gives the bounded-degree sofic presentation, Lemma~\ref{lem:m3-decoder} isolates the threshold arithmetic at $m=3$, Theorem~\ref{thm:block-invertibility} proves injectivity for every $m\ge 3$, Theorem~\ref{thm:m2-branch-locus} classifies the unique exceptional branch at $m=2$, and Theorem~\ref{thm:two-sided-conjugacy} upgrades the above-threshold picture to finite-memory conjugacy and finite-type closure. For background on sofic shifts and labeled graph presentations we refer to \cite{LindMarcus1995,Kitchens1998,Nasu1995}.

\begin{definition}[Sliding stabilized-window process]\label{def:Phi}
Define
\[
\Phi_m:\{0,1\}^{\mathbb Z}\longrightarrow (X_m)^{\mathbb Z}
\]
by
\[
(\Phi_m(a))_t:=\Fold_m(a_ta_{t+1}\cdots a_{t+m-1}),\qquad t\in\mathbb Z.
\]
Let
\[
Y_m:=\Phi_m(\{0,1\}^{\mathbb Z})\subset (X_m)^{\mathbb Z}.
\]
\end{definition}

\begin{definition}[The overlap graph]
Let $\mathcal G_m$ be the directed labeled graph defined as follows: the vertex set is
\[
V_m:=\{0,1\}^{m-1},
\]
and for each $u=(u_1,\dots,u_{m-1})\in V_m$ and each $c\in\{0,1\}$ there is an edge
\[
u \longrightarrow v=(u_2,\dots,u_{m-1},c)
\]
with label
\[
\ell(u,c):=\Fold_m(u_1,\dots,u_{m-1},c)\in X_m.
\]
\end{definition}

\begin{theorem}[Finite-state lift and bounded multiplicity]\label{thm:finite-state-lift}
For every $m\ge 2$:
\begin{enumerate}[label=(\roman*),leftmargin=*,itemsep=2pt]
  \item The shift $Y_m$ is exactly the label shift presented by $\mathcal G_m$. In particular, $Y_m$ is a sofic shift.
  \item The labeled graph $\mathcal G_m$ is resolving in both time directions: for each fixed vertex $u$, the two outgoing edges labeled by $c=0$ and $c=1$ have distinct labels, and for each fixed terminal vertex the two incoming edges also have distinct labels.
  \item The map $\Phi_m$ is finite-to-one, with degree at most $2^{m-1}$. More precisely, once the initial overlap state
  \[
  (a_0,\dots,a_{m-2})\in\{0,1\}^{m-1}
  \]
  is fixed, the entire raw sequence $a\in\{0,1\}^{\mathbb Z}$ is uniquely determined by its stabilized-window image $\Phi_m(a)$.
\end{enumerate}
\end{theorem}

\begin{proof}
(i) Let $a\in\{0,1\}^{\mathbb Z}$ and define vertices
\[
u_t:=(a_t,\dots,a_{t+m-2})\in\{0,1\}^{m-1}.
\]
Then
\[
u_t\to u_{t+1}
\]
is an edge in $\mathcal G_m$ labeled by
\[
\Fold_m(a_t,\dots,a_{t+m-1})=(\Phi_m(a))_t.
\]
Thus every point of $Y_m$ is realized as the label sequence of a bi-infinite path in $\mathcal G_m$.

Conversely, given any bi-infinite path $(u_t)_{t\in\mathbb Z}$ in $\mathcal G_m$, the overlap condition
\[
u_{t+1}=(u_{t,2},\dots,u_{t,m-1},c_t)
\]
defines a unique bi-infinite binary sequence $a$ whose length-$(m-1)$ windows are the $u_t$. The corresponding label sequence is exactly $\Phi_m(a)$. Hence $Y_m$ is the label shift of $\mathcal G_m$.

(ii) Fix $u\in V_m$. The two outgoing edges correspond to the length-$m$ words $u0$ and $u1$. Their Fibonacci values differ by exactly $F_{m+1}$:
\[
\Val_m(u1)-\Val_m(u0)=F_{m+1}.
\]
Since $0<F_{m+1}<F_{m+2}$, the two residues modulo $F_{m+2}$ are distinct. By Proposition~\ref{prop:finite-window-range}, distinct residues correspond to distinct labels in $X_m$. Therefore the two outgoing labels are distinct.

Now fix a terminal vertex $v=(v_1,\dots,v_{m-1})\in V_m$. Its two incoming edges come from the predecessor vertices
\[
(0,v_1,\dots,v_{m-2})
\qquad\text{and}\qquad
(1,v_1,\dots,v_{m-2}),
\]
both ending at $v$ with appended bit $v_{m-1}$. The corresponding length-$m$ words differ only in the first coordinate, so their Fibonacci values differ by exactly $F_2=1$. Hence their residues modulo $F_{m+2}$ are distinct, and the two incoming labels are distinct as well.

(iii) By (ii), once an initial vertex $u_0\in V_m$ and a label sequence $y=(y_t)_{t\in\mathbb Z}\in Y_m$ are given, there is at most one outgoing edge from the current vertex with the prescribed label $y_t$. This determines the forward path $(u_t)_{t\ge 0}$. The same item gives at most one incoming edge with prescribed label $y_{t-1}$ into a given terminal vertex $u_t$, so the path is uniquely determined backward as well. Hence the full bi-infinite path, and therefore the raw sequence $a$, are uniquely determined by $u_0$ and $y$. Since there are exactly $|V_m|=2^{m-1}$ possible initial vertices, the fiber cardinality of $\Phi_m$ over any $y\in Y_m$ is at most $2^{m-1}$.
\end{proof}

The arithmetic content of the threshold is already visible at $m=3$. Combining the three output residues $r_0,r_1,r_2$ isolates the sum $a_0+a_1$ together with the tail bit $a_4$; once $a_4$ is known, the remaining bits are recovered from the injective maps $(u,v)\mapsto u+2v$ on $\{0,1\}^2$. Figure~\ref{fig:m3-decoder-schematic} summarizes this reconstruction.

\begin{figure}[t]
\centering
\fbox{%
\parbox{0.9\textwidth}{%
\textbf{Step 1.} From $(r_0,r_1,r_2)$ compute $q=r_0-r_1-r_2 \pmod 5$.\\
\textbf{Step 2.} If $q\in\{0,1\}$ then $a_4=0$; if $q\in\{3,4\}$ then $a_4=1$; if $q=2$ use $(r_0,r_2)$ to separate the two branches.\\
\textbf{Step 3.} Recover $(a_2,a_3)$ from $t=r_2-3a_4\equiv a_2+2a_3 \pmod 5$.\\
\textbf{Step 4.} Recover $(a_0,a_1)$ from $u=r_0-3a_2\equiv a_0+2a_1 \pmod 5$.%
}}
\caption{Schematic of the three-window decoder. The quantity $q=r_0-r_1-r_2 \pmod 5$ isolates the tail bit except on a single branch, which is then resolved by $(r_0,r_2)$.}
\label{fig:m3-decoder-schematic}
\end{figure}

\begin{lemma}[Explicit three-window decoder]\label{lem:m3-decoder}
The map
\[
\Phi_{3,3}:\{0,1\}^{5}\longrightarrow X_3^3
\]
is injective.
\end{lemma}

\begin{proof}
Write the raw block as $a_0a_1a_2a_3a_4$ and let
\[
M:=F_5=5,
\qquad
r_t:=\Val_3\bigl(\Fold_3(a_ta_{t+1}a_{t+2})\bigr)\in\{0,1,2,3,4\},
\qquad 0\le t\le 2.
\]
Because $\Val_3|_{X_3}$ is a bijection, the stabilized output triple is equivalent to the residue triple $(r_0,r_1,r_2)$.

Using
\[
r_t\equiv a_t+2a_{t+1}+3a_{t+2}\pmod 5,
\]
set
\[
q:=r_0-r_1-r_2\pmod 5.
\]
The Fibonacci recurrence gives
\[
q\equiv a_0+a_1+2a_4\pmod 5.
\]
Since $a_0+a_1\in\{0,1,2\}$ and $2a_4\in\{0,2\}$, one has:
\[
q\in\{0,1\}\Longrightarrow a_4=0,
\qquad
q\in\{3,4\}\Longrightarrow a_4=1.
\]
Only the case $q=2$ requires an extra branch distinction. Then the two possible triples $(a_0,a_1,a_4)$ are
\[
(1,1,0)\qquad\text{or}\qquad (0,0,1).
\]
The corresponding diagnostic residues are summarized in Table~\ref{tab:m3-decoder}.

\begin{table}[t]
\centering
\caption{Diagnostic residues for the exceptional branch $q=2$ in Lemma~\ref{lem:m3-decoder}.}
\label{tab:m3-decoder}
\begin{tabular}{ccc}
\toprule
$(a_0,a_1,a_4)$ & possible $r_0$ & if $r_0=3$, possible $r_2$ \\
\midrule
$(1,1,0)$ & $\{1,3\}$ & $\{0,2\}$ \\
$(0,0,1)$ & $\{0,3\}$ & $\{1,4\}$ \\
\bottomrule
\end{tabular}
\end{table}

Indeed, in the first branch one has
\[
r_0\equiv 3+3a_2\pmod 5,
\qquad
r_2\equiv a_2+2a_3\pmod 5,
\]
whereas in the second branch
\[
r_0\equiv 3a_2\pmod 5,
\qquad
r_2\equiv a_2+2a_3+3\pmod 5.
\]
Hence $(q,r_0,r_2)$ determines $a_4$ uniquely.

Once $a_4$ is known, define
\[
t:=r_2-3a_4\pmod 5.
\]
Then
\[
t\equiv a_2+2a_3\pmod 5.
\]
Because the map
\[
\{0,1\}^2\to\{0,1,2,3\},\qquad (u,v)\mapsto u+2v
\]
is injective, $(a_2,a_3)$ is uniquely determined by $t$. Finally define
\[
u:=r_0-3a_2\pmod 5.
\]
Then
\[
u\equiv a_0+2a_1\pmod 5,
\]
and the same injectivity recovers $(a_0,a_1)$. Thus $a_0a_1a_2a_3a_4$ is uniquely determined by $(r_0,r_1,r_2)$, proving injectivity.
\end{proof}

\begin{theorem}[Sharp three-window threshold]\label{thm:block-invertibility}
Let $m\ge 3$. For each $n\ge m$, define
\[
\Phi_{m,n}:\{0,1\}^{n+m-1}\longrightarrow X_m^n
\]
by
\[
\Phi_{m,n}(a_0\ldots a_{n+m-2})
:=
\bigl(\Fold_m(a_t\ldots a_{t+m-1})\bigr)_{t=0}^{n-1}.
\]
Then $\Phi_{m,n}$ is injective.
\end{theorem}

\begin{proof}
If $m=3$, first consider $n=3$. This is exactly Lemma~\ref{lem:m3-decoder}. If $n>3$, the first three stabilized outputs determine the initial raw block
\[
a_0a_1a_2a_3a_4
\]
uniquely by Lemma~\ref{lem:m3-decoder}. In particular the initial overlap state
\[
(a_0,a_1)\in\{0,1\}^2
\]
is recovered. The right-resolving property in Theorem~\ref{thm:finite-state-lift}(ii) then recursively determines each subsequent raw bit from the current overlap state and the next stabilized symbol. Hence $\Phi_{3,n}$ is injective for every $n\ge 3$.

Assume from now on that $m\ge 4$.

Fix $a_0,\dots,a_{n+m-2}\in\{0,1\}$ and write
\[
M:=F_{m+2},
\qquad
r_t:=\Val_m\bigl(\Fold_m(a_t\ldots a_{t+m-1})\bigr)\in\{0,\dots,M-1\}.
\]
By Proposition~\ref{prop:finite-window-range}, the residue block $(r_t)_{t=0}^{n-1}$ is equivalent to the stabilized output block $\Phi_{m,n}(a_0\ldots a_{n+m-2})$.

For $0\le t\le m-3$, define
\[
q_t:=r_t-r_{t+1}-r_{t+2}\pmod M.
\]
Using
\[
r_t=\sum_{j=0}^{m-1} a_{t+j}F_{j+2},
\]
the Fibonacci recurrence gives
\[
q_t\equiv a_t+a_{t+1}+F_m a_{t+m+1}\pmod M.
\]
Indeed, the coefficients of $a_{t+j}$ cancel for $2\le j\le m-1$, the coefficient of $a_{t+m}$ is
\[
-F_{m+1}-F_m=-F_{m+2}\equiv 0\pmod M,
\]
and the coefficient of $a_{t+m+1}$ is
\[
-F_{m+1}\equiv F_m\pmod M.
\]

Since $m\ge 4$, one has $F_m\ge 3$. Hence the sets
\[
\{0,1,2\}
\qquad\text{and}\qquad
\{F_m,F_m+1,F_m+2\}
\]
are disjoint in $\{0,\dots,M-1\}$. Therefore each $q_t$ uniquely determines both
\[
a_{t+m+1}\in\{0,1\}
\qquad\text{and}\qquad
s_t:=a_t+a_{t+1}\in\{0,1,2\}.
\]
From $(q_0,\dots,q_{m-3})$ we recover all tail bits
\[
a_{m+1},a_{m+2},\dots,a_{2m-2},
\]
and all adjacent prefix sums
\[
s_0=a_0+a_1,\ s_1=a_1+a_2,\ \dots,\ s_{m-3}=a_{m-3}+a_{m-2}.
\]

Next, use the last two residues among the first $m$ windows. Define
\[
u:=r_{m-1}-\sum_{j=2}^{m-1} a_{m-1+j}F_{j+2}\pmod M.
\]
All terms except $a_{m-1}$ and $a_m$ have been removed, so
\[
u\equiv a_{m-1}F_2+a_mF_3=a_{m-1}+2a_m\pmod M.
\]
Because $a_{m-1}+2a_m\in\{0,1,2,3\}\subset [0,M-1]$, this congruence is an ordinary integer identity, and $(a_{m-1},a_m)$ is uniquely determined.

Similarly, define
\[
v:=r_{m-2}-\sum_{j=1}^{m-1} a_{m-2+j}F_{j+2}\pmod M.
\]
Then
\[
v\equiv a_{m-2}F_2=a_{m-2}\pmod M,
\]
so $a_{m-2}$ is uniquely determined.

Now recover the remaining prefix bits by backward recursion:
\[
a_{m-3}=s_{m-3}-a_{m-2},\quad
a_{m-4}=s_{m-4}-a_{m-3},\quad \dots,\quad
a_0=s_0-a_1.
\]
Thus the first $m$ stabilized outputs already determine the entire raw block
\[
a_0,\dots,a_{2m-2}.
\]

If $n=m$, the proof is complete. If $n>m$, the recovered initial overlap state
\[
(a_0,\dots,a_{m-2})\in\{0,1\}^{m-1}
\]
together with Theorem~\ref{thm:finite-state-lift}(ii) implies that each subsequent label uniquely determines the next raw bit. Iterating this right-resolving decoding recovers
\[
a_{2m-1},\dots,a_{n+m-2}.
\]
Therefore $\Phi_{m,n}$ is injective.
\end{proof}

\begin{corollary}[Explicit sliding-block decoder above threshold]\label{cor:explicit-decoder}
Assume $m\ge 3$, and let
\[
\Psi_m:\mathcal L_m(Y_m)\to\{0,1\}^{2m-1}
\]
be the inverse of $\Phi_{m,m}$ onto its image. Define
\[
\psi_m(w):=(\Psi_m(w))_1
\qquad\text{for }w\in\mathcal L_m(Y_m).
\]
Then for every $a\in\{0,1\}^{\mathbb Z}$ and $y=\Phi_m(a)$,
\[
a_t=\psi_m(y_t,\dots,y_{t+m-1})
\qquad\text{for all }t\in\mathbb Z.
\]
Thus the inverse factor map has memory $0$ and anticipation $m-1$.
\end{corollary}

\begin{proof}
For each $t\in\mathbb Z$, the block
\[
w_t:=(y_t,\dots,y_{t+m-1})
\]
lies in $\mathcal L_m(Y_m)$ and is the image under $\Phi_{m,m}$ of the raw block
\[
(a_t,\dots,a_{t+2m-2}).
\]
By Theorem~\ref{thm:block-invertibility}, $\Phi_{m,m}$ is injective, so
\[
\Psi_m(w_t)=(a_t,\dots,a_{t+2m-2}).
\]
Taking the first coordinate gives
\[
\psi_m(w_t)=a_t,
\]
as required.
\end{proof}

\begin{proposition}[Sharp low-window counterexamples]\label{prop:low-window-counterexamples}
The threshold in Theorem~\ref{thm:block-invertibility} is sharp.
\begin{enumerate}[label=(\roman*),leftmargin=*,itemsep=2pt]
  \item At $m=2$, sequence-level injectivity fails:
  \[
  \Phi_2(0^{\mathbb Z})=\Phi_2(1^{\mathbb Z})=(\ldots,00,00,00,\ldots).
  \]
  \item At $m=3$, the requirement $n\ge m$ is sharp:
  \[
  \Phi_{3,2}(0001)=\Phi_{3,2}(0110)=(000,001).
  \]
\end{enumerate}
\end{proposition}

\begin{proof}
For (i), $\Fold_2(00)=00$ and
\[
\Val_2(11)=F_2+F_3=3\equiv 0\pmod{F_4},
\]
so $\Fold_2(11)=00$ as well. Hence every length-$2$ window in both constant sequences stabilizes to $00$.

For (ii), one has
\[
\Fold_3(000)=000,\qquad \Fold_3(001)=001,
\]
and also
\[
\Val_3(011)=F_3+F_4=5\equiv 0\pmod{F_5},
\qquad
\Val_3(110)=F_2+F_3=3,
\]
so
\[
\Fold_3(011)=000,
\qquad
\Fold_3(110)=001.
\]
Thus the two raw blocks $0001$ and $0110$ have the same stabilized output pair.
\end{proof}

\begin{theorem}[Complete classification of the two-window branch locus]\label{thm:m2-branch-locus}
For $n\ge 1$, define
\[
\Phi_{2,n}:\{0,1\}^{n+1}\longrightarrow X_2^n,
\qquad
\Phi_{2,n}(a_0\ldots a_n):=\bigl(\Fold_2(a_ta_{t+1})\bigr)_{t=0}^{n-1},
\]
and write
\[
\mathcal Y_{2,n}:=\Phi_{2,n}(\{0,1\}^{n+1}).
\]
Then
\[
\Phi_{2,n}^{-1}\bigl((00)^n\bigr)=\{0^{n+1},1^{n+1}\},
\]
and for every $y\in\mathcal Y_{2,n}\setminus\{(00)^n\}$ one has
\[
|\Phi_{2,n}^{-1}(y)|=1.
\]
On two-sided path space,
\[
\Phi_2^{-1}\bigl((\ldots,00,00,\ldots)\bigr)=\{0^{\mathbb Z},1^{\mathbb Z}\},
\]
and every other point of $Y_2$ has singleton fiber.
\end{theorem}

\begin{proof}
At $m=2$ one has
\[
\Fold_2(00)=00,\qquad \Fold_2(11)=00,\qquad \Fold_2(01)=01,\qquad \Fold_2(10)=10.
\]
Thus the stabilized symbol $00$ records equality of adjacent raw bits, while $01$ and $10$ record the two possible local jumps.

Fix $y=(y_0,\dots,y_{n-1})\in\mathcal Y_{2,n}$. If $y=(00)^n$, then every adjacent pair in the raw block must consist of equal bits, so the only possibilities are $0^{n+1}$ and $1^{n+1}$. Hence
\[
\Phi_{2,n}^{-1}\bigl((00)^n\bigr)=\{0^{n+1},1^{n+1}\}.
\]

If $y\neq (00)^n$, choose $t$ with $y_t\in\{01,10\}$. Then $y_t$ determines the pair $(a_t,a_{t+1})$ exactly. For $s=t+1,\dots,n-1$, the symbol $y_s$ determines $a_{s+1}$ from $a_s$: if $y_s=00$ then $a_{s+1}=a_s$, if $y_s=01$ then $(a_s,a_{s+1})=(0,1)$, and if $y_s=10$ then $(a_s,a_{s+1})=(1,0)$. The same rule propagates uniquely backward for $s=t-1,\dots,0$. Therefore the entire raw block is uniquely determined, so every other fiber is a singleton.

The bi-infinite statement is proved in the same way. If $y=(\ldots,00,00,\ldots)$, then every adjacent raw pair is constant, so the only preimages are $0^{\mathbb Z}$ and $1^{\mathbb Z}$. Otherwise choose any $t$ with $y_t\in\{01,10\}$ and propagate the reconstruction uniquely in both time directions.
\end{proof}

\begin{corollary}[Exact entropy and KL defect formulas at the two-window threshold]\label{cor:m2-defect-formulas}
Fix $n\ge 1$, let $A\sim P$ be a law on $\{0,1\}^{n+1}$, and set $Y:=\Phi_{2,n}(A)$. Write
\[
p_0:=P(0^{n+1}),
\qquad
p_1:=P(1^{n+1}),
\qquad
q:=p_0+p_1,
\]
and define the binary entropy
\[
h(r):=-r\log r-(1-r)\log(1-r),
\qquad 0\le r\le 1.
\]
Then
\[
H_P(A\mid Y)=q\,h\Bigl(\frac{p_0}{q}\Bigr)\le q\log 2,
\]
with the convention that the right-hand side is $0$ when $q=0$.

If $Q$ is another law on $\{0,1\}^{n+1}$, with
\[
q_0:=Q(0^{n+1}),
\qquad
q_1:=Q(1^{n+1}),
\qquad
s:=q_0+q_1,
\]
then the KL defect vanishes when $q=0$. If $q>0$ and $s>0$, one has
\begin{equation}\label{eq:m2-kl-defect}
\begin{aligned}
D_{\KL}(P\|Q)&-D_{\KL}\bigl((\Phi_{2,n})_*P\|(\Phi_{2,n})_*Q\bigr)
\\
&=
q\Biggl[
\frac{p_0}{q}\log\frac{p_0/q}{q_0/s}
+
\frac{p_1}{q}\log\frac{p_1/q}{q_1/s}
\Biggr].
\end{aligned}
\end{equation}
Here terms with $p_i=0$ are interpreted as $0$.
\end{corollary}

\begin{proof}
By Theorem~\ref{thm:m2-branch-locus}, every conditional fiber of $\Phi_{2,n}$ is a singleton except the branch fiber over $(00)^n$, which is exactly $\{0^{n+1},1^{n+1}\}$. Therefore $H_P(A\mid Y=y)=0$ for $y\neq (00)^n$, while on the branch fiber the conditional law is
\[
P(A=0^{n+1}\mid Y=(00)^n)=\frac{p_0}{q},
\qquad
P(A=1^{n+1}\mid Y=(00)^n)=\frac{p_1}{q}.
\]
Hence
\[
H_P(A\mid Y)=P(Y=(00)^n)\,h\Bigl(\frac{p_0}{q}\Bigr)=q\,h\Bigl(\frac{p_0}{q}\Bigr),
\]
and the bound follows from $h(r)\le \log 2$.

For the KL identity, all singleton fibers cancel when one compares the raw divergence with the pushed-forward divergence. Only the branch fiber contributes:
\begin{equation*}
\begin{aligned}
D_{\KL}(P\|Q)&-D_{\KL}\bigl((\Phi_{2,n})_*P\|(\Phi_{2,n})_*Q\bigr)
\\
&=
p_0\log\frac{p_0}{q_0}
+
p_1\log\frac{p_1}{q_1}
-
q\log\frac{q}{s}.
\end{aligned}
\end{equation*}
If $q=0$ this is $0$. Otherwise $s>0$ by assumption, so factor out $q$ and rewrite the remaining expression in terms of the two-point distributions on the branch fiber, which gives \eqref{eq:m2-kl-defect}.
\end{proof}

\begin{corollary}[Asymptotic sufficiency off the branch fiber]\label{cor:m2-asymptotic-sufficiency}
Let $\{P_{u,n}:u\in\mathcal U\}$ be any family of laws on $\{0,1\}^{n+1}$, and suppose that
\[
\sup_{u\in\mathcal U} P_{u,n}\bigl(\{0^{n+1},1^{n+1}\}\bigr)\le \varepsilon_n.
\]
Then
\[
\sup_{u\in\mathcal U} H_{P_{u,n}}\bigl(A\mid \Phi_{2,n}(A)\bigr)\le \varepsilon_n\log 2.
\]
In particular, if $\varepsilon_n$ decays exponentially, then two-window stabilized readout is exponentially asymptotically sufficient for this family.
\end{corollary}

\begin{proof}
Apply Corollary~\ref{cor:m2-defect-formulas} and use $h(r)\le \log 2$.
\end{proof}

\begin{corollary}[Mod-zero invertibility at the two-window threshold]\label{cor:m2-mod-zero}
Let $\nu$ be a shift-invariant Borel probability measure on $\{0,1\}^{\mathbb Z}$ such that
\[
\nu(\{0^{\mathbb Z},1^{\mathbb Z}\})=0,
\]
and set $\eta:=(\Phi_2)_*\nu$. Then
\[
\Phi_2:\bigl(\{0,1\}^{\mathbb Z},\nu\bigr)\longrightarrow (Y_2,\eta)
\]
is a measure-theoretic isomorphism.
\end{corollary}

\begin{proof}
By Theorem~\ref{thm:m2-branch-locus}, the restriction of $\Phi_2$ to
\[
\{0,1\}^{\mathbb Z}\setminus\{0^{\mathbb Z},1^{\mathbb Z}\}
\]
is one-to-one, with image
\[
Y_2\setminus\{(\ldots,00,00,\ldots)\}.
\]
Since $\Phi_2$ is continuous, this restricted map is Borel, and its inverse on the image is Borel by the standard Lusin--Souslin theorem for injective Borel maps between Polish spaces. Both the branch set and its image have measure zero under $\nu$ and $\eta$, respectively. Hence the restriction of $\Phi_2$ furnishes an invertible, shift-equivariant correspondence modulo null sets, which is exactly a measure-theoretic isomorphism.
\end{proof}

\begin{corollary}[Full block complexity of the one-sided stabilized image]\label{cor:one-sided-complexity}
Let
\[
p_{Y_m^+}(n):=\#\mathcal L_n(Y_m^+)
\]
denote the length-$n$ word complexity of the one-sided stabilized image
\[
Y_m^+:=\Phi_m^+(\{0,1\}^{\mathbb N}),
\qquad
(\Phi_m^+(a))_t:=\Fold_m(a_t\ldots a_{t+m-1}).
\]
If $m\ge 3$ and $n\ge m$, then
\[
p_{Y_m^+}(n)=2^{n+m-1}.
\]
\end{corollary}

\begin{proof}
The domain of $\Phi_{m,n}$ has cardinality $2^{n+m-1}$. By Theorem~\ref{thm:block-invertibility}, $\Phi_{m,n}$ is injective, so its image has the same cardinality. But that image is precisely the set of length-$n$ blocks appearing in $Y_m^+$.
\end{proof}

\begin{corollary}[One-sided stabilized-window conjugacy]\label{cor:one-sided-conjugacy}
If $m\ge 3$, the one-sided sliding stabilized-window map
\[
\Phi_m^+:\bigl(\{0,1\}^{\mathbb N},\sigma\bigr)\longrightarrow (Y_m^+,\sigma)
\]
is a topological conjugacy.
\end{corollary}

\begin{proof}
If $\Phi_m^+(a)=\Phi_m^+(b)$ for $a,b\in\{0,1\}^{\mathbb N}$, then for every $n\ge m$ the two sequences have the same image under $\Phi_{m,n}$. By Theorem~\ref{thm:block-invertibility}, their first $n+m-1$ raw bits coincide. Letting $n\to\infty$ shows $a=b$. Thus $\Phi_m^+$ is injective.

It is continuous, shift-equivariant, and surjective onto $Y_m^+$ by definition. Since $\{0,1\}^{\mathbb N}$ is compact and $Y_m^+$ is Hausdorff, a continuous bijection onto $Y_m^+$ is a homeomorphism. Hence $\Phi_m^+$ is a topological conjugacy.
\end{proof}

\begin{theorem}[Two-sided finite-memory conjugacy and finite-type closure]\label{thm:two-sided-conjugacy}
Assume $m\ge 3$. Then:
\begin{enumerate}[label=(\roman*),leftmargin=*,itemsep=2pt]
  \item The two-sided stabilized-window map
  \[
  \Phi_m:\bigl(\{0,1\}^{\mathbb Z},\sigma\bigr)\to (Y_m,\sigma)
  \]
  is injective, hence a topological conjugacy onto its image.
  \item There exists a local decoder
  \[
  \psi_m:\mathcal L_m(Y_m)\to\{0,1\}
  \]
  such that for every $a\in\{0,1\}^{\mathbb Z}$ and $y=\Phi_m(a)$,
  \[
  a_t=\psi_m(y_t,\dots,y_{t+m-1})
  \qquad\text{for all }t\in\mathbb Z.
  \]
  Thus the inverse conjugacy is a sliding-block code with memory $0$ and anticipation $m-1$.
  \item The image shift $Y_m$ is a subshift of finite type of memory at most $m$.
\end{enumerate}
\end{theorem}

\begin{proof}
(ii) is exactly Corollary~\ref{cor:explicit-decoder}.

If $\Phi_m(a)=\Phi_m(b)$, then for every $t$ the corresponding length-$m$ output blocks agree, so by (ii) one has
\[
a_t=\psi_m(y_t,\dots,y_{t+m-1})=b_t.
\]
Hence $a=b$, proving injectivity. Since $\Phi_m$ is continuous, shift-equivariant, and surjective onto $Y_m$ by definition, (i) follows.

For (iii), let $y\in (X_m)^{\mathbb Z}$. For each $t$, if the length-$m$ block
\[
w_t:=(y_t,\dots,y_{t+m-1})
\]
belongs to $\mathcal L_m(Y_m)$, define
\[
B_t:=\Psi_m(w_t)\in\{0,1\}^{2m-1}.
\]
If $y=\Phi_m(a)$ for some $a$, then these decoded blocks must satisfy the overlap consistency condition
\begin{equation}\label{eq:decoded-overlap}
(B_t)_2(B_t)_3\cdots (B_t)_{2m-1}
=
(B_{t+1})_1(B_{t+1})_2\cdots (B_{t+1})_{2m-2}
\qquad\text{for all }t\in\mathbb Z.
\end{equation}
Conversely, suppose that every length-$m$ block of $y$ lies in $\mathcal L_m(Y_m)$ and that the decoded blocks satisfy \eqref{eq:decoded-overlap} for all $t$. Then the blocks $(B_t)$ glue uniquely to a bi-infinite raw sequence $a\in\{0,1\}^{\mathbb Z}$, and by construction every length-$m$ output block of $\Phi_m(a)$ agrees with the corresponding block of $y$. Hence $\Phi_m(a)=y$.

The conditions ``$w_t\in\mathcal L_m(Y_m)$'' and \eqref{eq:decoded-overlap} depend only on the length-$(m+1)$ block $(y_t,\dots,y_{t+m})$. Since there are only finitely many such blocks, they define a finite forbidden list. Therefore $Y_m$ is a subshift of finite type of memory at most $m$.
\end{proof}

\begin{remark}[Decoder complexity]\label{rem:decoder-complexity}
The inverse code $\psi_m$ of Theorem~\ref{thm:two-sided-conjugacy}(ii) reads a sliding window of length $m$ and outputs one raw bit. Equivalently, one may implement the decoder by maintaining an overlap state of size $m-1$. The arithmetic cost is $O(m)$ per output symbol, since each decoding step involves at most $m$ Fibonacci-weighted additions modulo $F_{m+2}$.
\end{remark}

\begin{proposition}[Immediate dynamical corollaries above threshold]\label{prop:conjugacy-consequences}
Assume $m\ge 3$.
\begin{enumerate}[label=(\roman*),leftmargin=*,itemsep=2pt]
  \item For every $n\ge 1$,
  \[
  \#\operatorname{Fix}(\sigma^n|_{Y_m})=2^n,
  \qquad
  \zeta_{Y_m}(z)=\frac{1}{1-2z}.
  \]
  \item Pushforward by $\Phi_m$ induces a bijection between the shift-invariant Borel probability measures on $\{0,1\}^{\mathbb Z}$ and those on $Y_m$.
  \item Ergodicity, weak mixing, mixing, Bernoulli-ness, $K$-property, and measure-theoretic entropy are transported exactly through the conjugacy.
\end{enumerate}
\end{proposition}

\begin{proof}
All three items are immediate from Theorem~\ref{thm:two-sided-conjugacy}. Fixed-point counts and the Artin--Mazur zeta function are conjugacy invariants. The inverse sliding-block code gives the bijection of invariant measures. The dynamical statements in (iii) are standard invariants of topological or measure-theoretic conjugacy; see \cite{Walters1982}. The thermodynamic consequences for continuous potentials are recorded separately in Theorem~\ref{thm:continuous-potentials}.
\end{proof}

\begin{remark}\label{rem:exact-recoding}
For $m\ge 3$ the stabilization image is not merely a sofic factor of bounded degree. It is a subshift of finite type topologically conjugate to the full two-shift by a finite-memory code, and the only obstruction below threshold is the explicit branch locus of Theorem~\ref{thm:m2-branch-locus}.
\end{remark}

\begin{remark}\label{rem:sofic-comparison}
In the terminology of \cite{LindMarcus1995}, $\mathcal{G}_m$ is a right-resolving presentation of the sofic shift $Y_m$, and $\Phi_m$ has degree at most $2^{m-1}$ (Theorem~\ref{thm:finite-state-lift}(iii)). A right-resolving factor of degree~1 between irreducible shifts of finite type is a conjugacy \cite[Theorem~8.1.16]{LindMarcus1995}; see also \cite{Ashley1988} for the equal-entropy resolving theory and \cite{Boyle1983} for the lower-entropy obstruction. The content of the threshold theorem is that degree~1 is achieved at $m=3$ and not before; this is an arithmetic fact that cannot be read off from the general sofic framework and requires the explicit decoder of Lemma~\ref{lem:m3-decoder}.
\end{remark}

\begin{proposition}[Measure-theoretic entropy preservation under the stabilized-window factor]\label{prop:measure-entropy-preservation}
Let $\nu$ be a shift-invariant Borel probability measure on $\{0,1\}^{\mathbb Z}$, and let
\[
\eta:=(\Phi_m)_*\nu
\]
be its image on $Y_m\subset (X_m)^{\mathbb Z}$. Then
\[
h_\eta(\sigma)=h_\nu(\sigma).
\]
\end{proposition}

\begin{proof}
For $n\ge 1$, let
\[
A_{0:n+m-2}:=(A_0,\dots,A_{n+m-2})
\]
denote the raw binary block of length $n+m-1$, and let
\[
Y_{0:n-1}:=(Y_0,\dots,Y_{n-1}),\qquad Y_t=\Fold_m(A_t\cdots A_{t+m-1}),
\]
denote the stabilized block of length $n$.

Since $Y_{0:n-1}$ is a deterministic function of $A_{0:n+m-2}$, one has
\[
H_\eta(Y_{0:n-1})\le H_\nu(A_{0:n+m-2}).
\]
Conversely, by Theorem~\ref{thm:finite-state-lift}(iii), once the initial overlap state
\[
(A_0,\dots,A_{m-2})\in\{0,1\}^{m-1}
\]
is specified, the entire raw block $A_{0:n+m-2}$ is uniquely determined by the stabilized block $Y_{0:n-1}$. Therefore
\[
H_\nu(A_{0:n+m-2}\mid Y_{0:n-1})\le (m-1)\log 2.
\]
By the chain rule,
\[
H_\nu(A_{0:n+m-2})
=
H_\eta(Y_{0:n-1})
+
H_\nu(A_{0:n+m-2}\mid Y_{0:n-1}),
\]
hence
\[
0\le H_\nu(A_{0:n+m-2})-H_\eta(Y_{0:n-1})\le (m-1)\log 2.
\]
Divide by $n$ and let $n\to\infty$. The additive discrepancy is $O(1)$, so it disappears after normalization. The left-hand side converges to $h_\nu(\sigma)$ because shifting the block length by the fixed amount $m-1$ does not affect the entropy rate, while the right-hand side converges to $h_\eta(\sigma)$. Thus
\[
h_\eta(\sigma)=h_\nu(\sigma).
\]
\end{proof}

\begin{corollary}[Sequence-level entropy preservation]\label{cor:entropy-preservation}
For every $m\ge 2$,
\[
h_{\mathrm{top}}(Y_m)=\log 2.
\]
\end{corollary}

\begin{proof}
The map $\Phi_m$ is a continuous factor map from the full two-shift onto $Y_m$, so
\[
h_{\mathrm{top}}(Y_m)\le h_{\mathrm{top}}(\{0,1\}^{\mathbb Z})=\log 2.
\]
\[
\text{Let }\nu_{1/2}:=\Bigl(\tfrac12,\tfrac12\Bigr)^{\otimes \mathbb Z}
\]
be the Bernoulli measure on the full two-shift, and let
\[
\eta_{1/2}:=(\Phi_m)_*\nu_{1/2}.
\]
By Proposition~\ref{prop:measure-entropy-preservation},
\[
h_{\eta_{1/2}}(\sigma)=h_{\nu_{1/2}}(\sigma)=\log 2.
\]
Since $\eta_{1/2}$ is a shift-invariant Borel probability measure on $Y_m$, the variational principle \cite{Walters1982} gives
\[
h_{\mathrm{top}}(Y_m)\ge \log 2.
\]
Therefore equality holds.
\end{proof}

\begin{remark}
For every $m\ge 2$ the stabilized-window factor preserves entropy. Once $m\ge 3$, the picture is much more rigid: $Y_m$ is conjugate to the full two-shift, so periodic data, invariant measures, pressure, and equilibrium states are transported exactly. Proposition~\ref{prop:low-window-counterexamples} shows that this threshold is sharp.
\end{remark}


